\documentclass{xStandalone}

\begin{document}
\begin{tikzpicture}

\draw[ultra thin,gray] (-0.5,-2.8) rectangle (8.9,2.8);
    
\draw (-0.3,-1.3) coordinate (A) rectangle (+0.3,+1.3) coordinate (B);

\path ($(A)!0.5!(B)$) coordinate(O) ++(0.02,0) node[text width=1em] {\small \hspace{0.13em}X射线管};

\draw[-latex] (O-|B) -- ++(1.5,0) coordinate (L) -- ++(2,0) coordinate (S);

\draw[fill=gray] ($(L)+(-0.3,+0.1)$) rectangle ($(L)+(-0.1,+0.6)$);
\draw[fill=gray] ($(L)+(+0.3,+0.1)$) rectangle ($(L)+(+0.1,+0.6)$);
\draw[fill=gray] ($(L)+(-0.3,-0.1)$) rectangle ($(L)+(-0.1,-0.6)$);
\draw[fill=gray] ($(L)+(+0.3,-0.1)$) rectangle ($(L)+(+0.1,-0.6)$);

\path (L) node[above=0.6cm] {\small 光阑};

\draw[fill=black!70!white] ($(S)+(0,-0.4)$) coordinate (S1) rectangle ($(S)+(0.8,+0.4)$) coordinate(S2);

\path ($(S1)!0.5!(S2)$) coordinate (Sc) node[above=0.6cm] {\small 散射物质} node[below=0.5cm] {\scriptsize （石墨）};

\draw[gray,very thick] (Sc-|S2) coordinate (Sx) ++(60:3) arc[radius=3,start angle=60,end angle=-60];

\draw[-latex] (Sx) -- node[below] {\small 散射光线} ++(0:3) coordinate (Sx1) node[right] {\small 光谱仪};
\draw[-latex] (Sx) -- ++(35:3) coordinate (Sx2);

\draw pic["{$\phi$}",
draw=black,
angle radius=0.4cm,
angle eccentricity=1.8] 
{angle=Sx1--Sx--Sx2};

\end{tikzpicture}
\end{document}